\newpage
\phantomsection
\label{prf:rational-perturbation-lemma}

\begin{remark*}[Return]
\hyperref[lem:rational-perturbation-lemma]{Return to Theorem}
\end{remark*}

\begin{theorem*}[Rational Perturbation Lemma]
If $q\in\mathbb{Q}$ and $\varepsilon>0$, then there exists $r\in\mathbb{Q}$
such that
\[
0<|r-q|<\varepsilon.
\]
\end{theorem*}

\begin{proof}
Let $q\in\mathbb{Q}$ and let $\varepsilon>0$. By
\hyperref[lem:rational-upper-lower-perturbation]{Lemma~\ref*{lem:rational-upper-lower-perturbation}},
there exist rational numbers $r,s\in\mathbb{Q}$ such that
\[
q-\varepsilon<r<q<s<q+\varepsilon.
\]
We may use the point $r$. Since $r<q$, we have
\[
0<q-r<\varepsilon.
\]
Because $r-q<0$, the definition of absolute value gives
\[
|r-q|=q-r.
\]
Therefore
\[
0<|r-q|<\varepsilon.
\]
Since $r\in\mathbb{Q}$, this proves the claim.
\end{proof}

\begin{remark*}[Proof structure]
The upper/lower perturbation lemma already gives rational points strictly on
both sides of $q$ within distance $\varepsilon$. To obtain the present result,
it is enough to choose one of those points, say $r<q$, and rewrite the
inequality
\[
q-\varepsilon<r<q
\]
as
\[
0<q-r<\varepsilon.
\]
Since $r-q<0$, one has $|r-q|=q-r$, which yields the desired estimate.
\end{remark*}

\begin{dependencies}
\begin{itemize}
  \item \hyperref[lem:rational-upper-lower-perturbation]{Lemma~\ref*{lem:rational-upper-lower-perturbation}}
\end{itemize}
\end{dependencies}

